\section{Periodic Point Detection (Feature Finder)}
\label{sec:feature-finder}

Given a viewer location and a search radius, \FractalShark{} can automatically
locate the center of a nearby mini-Mandelbrot copy---i.e.\ the unique parameter
$c$ at which the critical orbit is exactly periodic.  This capability is
implemented by the \code{FeatureFinder} class (files
\code{FeatureFinder.h}/\code{FeatureFinder.cpp}), whose design is heavily
influenced by the periodic-point finder in the Imagina fractal viewer \cite{Imagina}.  The
algorithm proceeds in two phases:
\begin{enumerate}
\item \textbf{Phase~A (T-space).}
  Detect the candidate period $p$ and perform a coarse Newton iteration using
  the renderer's native floating-point type~$T$ (which may be \code{float},
  \code{double}, or \code{HDRFloat}).
\item \textbf{Phase~B (MPF refinement).}
  Polish the candidate to arbitrary precision using GMP \code{mpf\_t}
  arithmetic, employing a mixed Newton--Halley iteration with precision
  management modeled after Imagina.
\end{enumerate}
Three evaluation back-ends are available for Phase~A---direct iteration,
perturbation theory (PT), and linear approximation (LA)---so that the finder
can operate at any zoom depth supported by the renderer.

\subsection{Motivation: mini-Mandelbrot copies}
\label{subsec:ff-motivation}

The Mandelbrot set contains infinitely many small-scale copies of itself,
commonly called \emph{mini-brots}.  Each copy is the nucleus of a hyperbolic
component of period~$p$: a connected region of the parameter plane in which the
critical orbit $z_0=0,\;z_{n+1}=z_n^2+c$ converges to an attracting cycle of
exact period~$p$.  The center of the component---the \emph{nucleus}---is the
unique parameter $c^*$ satisfying
\begin{equation}\label{eq:ff-nucleus}
z_p(c^*) = 0,\qquad z_0 = 0,
\end{equation}
where $z_p$ denotes the $p$-fold iterate of the critical point.

Locating these nuclei is useful for several reasons:
\begin{itemize}
\item \textbf{Automatic navigation.}  Given a user click, the finder can snap
  to the nearest mini-brot center and report its period and intrinsic size,
  enabling one-click deep-zoom targeting.
\item \textbf{Scale estimation.}  The derivative $\partial z_p/\partial c$ at
  the nucleus determines the local magnification scale of the mini-brot
  relative to the full set, which informs precision requirements for rendering.
\end{itemize}
At extreme zoom depths the nucleus coordinates may require thousands of bits of
precision, so the finder must combine fast low-precision search with
high-precision refinement.

\subsection{Design and implementation}
\label{subsec:ff-design}

\paragraph{Class template.}
\code{FeatureFinder<IterType, T, PExtras>} is a final class template
parameterized by the iteration-count type (\code{uint32\_t} or
\code{uint64\_t}), the working floating-point type~$T$, and a
\code{PerturbExtras} tag that selects compression behavior.  It inherits
convenience typedefs from \code{TemplateHelpers} and defines a nested
\code{Params} struct with tunable constants (maximum Newton iterations,
relative step tolerances $2^{-40}$ and $2^{-80}$, an optional absolute
residual accept threshold, and a print flag).

\paragraph{Evaluator policies (strategy pattern).}
Phase~A must iterate the critical orbit to detect the period and compute
derivatives.  Three evaluator structs encapsulate this:
\begin{itemize}
\item \code{DirectEvaluator} --- iterates $z\leftarrow z^2+c$ from scratch,
  suitable when no reference orbit is available.
\item \code{PTEvaluator} --- uses a precomputed reference orbit and
  perturbation delta $\Delta c$ to evaluate the orbit, decompressing
  reference data on the fly via \code{RuntimeDecompressor}.
\item \code{LAEvaluator} --- uses the linear-approximation (LA) coefficients
  from \code{LAReference} to skip early iterations, falling back to
  perturbation for the remaining tail.
\end{itemize}
Each evaluator exposes a single templated method \code{Eval<FindPeriod>}.
When \code{FindPeriod=true}, it searches for the smallest $n$ at which a
periodicity trigger fires; when \code{FindPeriod=false}, it evaluates at a
known period and returns the residual, $\partial z/\partial c$, and
$z_{\mathrm{coeff}}$.

\paragraph{Common Newton loop (\code{FindPeriodicPoint\_Common}).}
All three \code{FindPeriodicPoint} overloads delegate to this template,
which:
\begin{enumerate}
\item Calls the evaluator with \code{FindPeriod=true} to obtain a candidate
  period~$p$.
\item Applies an initial Newton correction
  $c \leftarrow c - z_p / (\partial z_p/\partial c)$ in high precision
  (the canonical $c$ is maintained as a \code{HighPrecision} value and
  projected to $T$ for each evaluation).
\item Runs up to \code{MaxNewtonIters} Newton steps in $T$-space, stopping when
  either
  $|\Delta c|^2 \le |c|^2 \cdot (2^{-40})^2$ (relative step) or
  $|z_p|^2 \le |c|^2 \cdot (2^{-40})^2$ (relative residual).
\item Performs a final correction pass and stores the result as a
  \code{PeriodicPointCandidate} inside \code{FeatureSummary}, along with the
  MPF precision needed for Phase~B.
\end{enumerate}

\paragraph{Phase~B: MPF refinement (\code{RefinePeriodicPoint\_HighPrecision}).}
This method retrieves the candidate from \code{FeatureSummary}, converts
its coordinates to GMP \code{mpf\_t} at the candidate's stored precision,
and calls the Imagina-style mixed Newton--Halley polisher described in
\cref{subsec:ff-newton,subsec:ff-halley,subsec:ff-mixed-prec}.  On
completion the refined coordinates are written back to \code{FeatureSummary}
and the candidate is marked as refined.

\paragraph{Parallel orbit evaluation.}
The high-precision forward evaluation
(\code{Evaluate\-Critical\-Orbit\-And\-Derivs}) must perform several large
\code{mpf\_mul} operations per iteration.  To overlap these, the code spawns
seven spin-based worker threads (four for derivative-precision products,
three for coordinate-precision products).  Workers spin on per-slot
\code{std::atomic\allowbreak<uint64\_t>} generation counters;
the main thread
publishes operand pointers, bumps the generation, and spin-waits on a
matching done counter.  Each slot is cache-line-aligned (\code{alignas(64)})
to avoid false sharing.

\subsection{Newton--Raphson iteration}
\label{subsec:ff-newton}

\paragraph{Root-finding formulation.}
Finding a period-$p$ nucleus amounts to solving $F(c)=0$ where
\begin{equation}\label{eq:ff-Fc}
F(c) \;=\; z_p(c),\qquad
z_0 = 0,\quad z_{n+1} = z_n^2 + c.
\end{equation}
The algorithm decomposes the root-finding problem \eqref{eq:ff-Fc} into two nested loops: an \emph{outer loop} that
updates the candidate~$c$ via Newton's method, and an \emph{inner loop} that
evaluates $F(c)$ and its derivatives by iterating the critical orbit for
$p$~steps.  The outer loop is inherently sequential---each Newton step
depends on the previous one---but the inner loop is the computational
bottleneck and is amenable to GPU acceleration
(\cref{subsec:ff-inner-gpu}).

\subsubsection{Outer loop: Newton step}
\label{subsubsec:ff-outer}

Given the outputs of the inner loop---$z_p$, $\partial z_p/\partial c$, and
optionally $\partial^2 z_p/\partial c^2$---Newton's method produces the
update
\begin{equation}\label{eq:ff-newton-step}
\Delta c \;=\; \frac{F(c)}{F'(c)}
         \;=\; \frac{z_p}{\partial z_p / \partial c},
\qquad c \leftarrow c - \Delta c.
\end{equation}
Writing $z_p = z_x + i\,z_y$ and $\partial z_p/\partial c = d_x + i\,d_y$,
the complex division expands to:
\begin{align}
\Delta c_x &= \frac{z_x\,d_x + z_y\,d_y}{d_x^2 + d_y^2}, \\[4pt]
\Delta c_y &= \frac{z_y\,d_x - z_x\,d_y}{d_x^2 + d_y^2}.
\end{align}
(When the Halley gate is satisfied, a higher-order step using
$\partial^2 z_p/\partial c^2$ replaces the Newton step; see
\cref{subsec:ff-halley}.)

\paragraph{Convergence.}
Standard Newton--Raphson converges quadratically near a simple root.  The
Phase~A loop in \code{FindPeriodicPoint\_Common} uses two independent
stopping criteria, whichever triggers first:
\begin{enumerate}
\item \emph{Relative step:}
  $|\Delta c|^2 \le |c|^2 \cdot \epsilon_{\mathrm{step}}^2$
  with $\epsilon_{\mathrm{step}} = 2^{-40}$.
\item \emph{Relative residual:}
  $|z_p|^2 \le |c|^2 \cdot \epsilon_{\mathrm{diff}}^2$
  with $\epsilon_{\mathrm{diff}} = 2^{-40}$.
\end{enumerate}
An optional absolute residual accept (\code{Eps2Accept}) provides a third
early-out that can be enabled via \code{Params}.

\paragraph{Period detection trigger.}
Before Newton iteration begins, the evaluator must identify the candidate
period.  The criterion used is: at iteration~$n$, if
\begin{equation}\label{eq:ff-period-trigger}
|z_n|^2 \;<\; R^2 \cdot \left|\frac{\partial z_n}{\partial c}\right|^2,
\end{equation}
then $p = n$ is accepted as the candidate period, where $R$ is the search
radius.  This test detects the moment the orbit passes close to the origin
relative to the local stretching factor $|\partial z/\partial c|$, which is
characteristic of near-periodic behavior.  A dynamic tightening mechanism
(implemented in the \code{PeriodicityPP} helper) progressively shrinks the
effective $R^2$ during the scan to reduce false triggers.

\subsubsection{Inner loop: forward evaluation}
\label{subsubsec:ff-inner}

\begin{figure}[!htbp]
\centering
\begin{tikzpicture}[
    >=Stealth,
    box/.style={draw, rounded corners, minimum width=2.8cm, minimum height=0.7cm,
                align=center, font=\small},
    outerbox/.style={box, fill=blue!8},
    innerbox/.style={box, fill=orange!10},
    precbox/.style={font=\scriptsize\itshape, text=gray},
    every edge/.style={draw, ->},
]
  % Outer loop
  \node[outerbox] (init) {Initial guess $c$};
  \node[outerbox, below=0.5cm of init] (inner) {Inner loop\\(evaluate $z_p$, $\mathit{dzdc}_p$, $d2_p$)};
  \node[outerbox, below=0.5cm of inner] (step) {Newton step\\$c \leftarrow c - z_p/\mathit{dzdc}_p$};
  \node[outerbox, below=0.5cm of step] (conv) {Converged?};
  \node[outerbox, below=0.5cm of conv] (done) {Return $c^*$};

  \draw[->] (init) -- (inner);
  \draw[->] (inner) -- (step);
  \draw[->] (step) -- (conv);
  \draw[->] (conv) -- node[right, font=\scriptsize] {yes} (done);
  \draw[->] (conv.west) -- ++(-1.5,0) |- node[near start, left, font=\scriptsize] {no} (inner.west);

  % Inner loop detail
  \node[innerbox, right=2.5cm of inner, yshift=1.0cm] (orbit) {$z \leftarrow z^2 + c$};
  \node[precbox, right=0.1cm of orbit.east] {coord};
  \node[innerbox, below=0.3cm of orbit] (dzdc) {$\mathit{dzdc} \leftarrow 2z \cdot \mathit{dzdc} + 1$};
  \node[precbox, right=0.1cm of dzdc.east] {deriv};
  \node[innerbox, below=0.3cm of dzdc] (d2) {$d2 \leftarrow 2(\mathit{dzdc}^2 + z \cdot d2)$};
  \node[precbox, right=0.1cm of d2.east] {${\sim}53$-bit};

  \node[font=\scriptsize, above=0.15cm of orbit] {Repeat $p$ times:};

  \begin{scope}[on background layer]
    \node[draw, dashed, rounded corners, fill=orange!4,
          fit=(orbit)(dzdc)(d2), inner xsep=6pt, inner ysep=10pt,
          label={[font=\scriptsize\bfseries]above:Inner loop (GPU-amenable)}] (innerfit) {};
  \end{scope}

  \draw[->, dashed] (inner.east) -- (innerfit.west);

  % Outer loop label
  \begin{scope}[on background layer]
    \node[draw, dotted, rounded corners, fill=blue!3,
          fit=(init)(done)(conv), inner xsep=8pt, inner ysep=8pt,
          label={[font=\scriptsize\bfseries]above:Outer loop (CPU)}] {};
  \end{scope}
\end{tikzpicture}
\caption[Newton--Raphson two-level loop structure]{Newton--Raphson two-level loop structure.  The outer loop
  (left) updates the candidate~$c$ via Newton steps on the CPU.  Each
  step calls the inner loop (right), which iterates $p$~times to
  evaluate $z_p$, $\mathit{dzdc}_p$, and $d2_p$ at three different
  precision levels.  The inner loop is the computational bottleneck and
  the target for GPU acceleration.}
\label{fig:newton-loops}
\end{figure}

Each Newton step requires evaluating $F(c) = z_p(c)$ and its first two
derivatives at the current guess~$c$, as shown in \cref{fig:newton-loops}.
This evaluation is done by iterating the
critical orbit from $z_0 = 0$ for exactly $p$~steps, accumulating three
quantities simultaneously at \emph{different precision levels}:

\begin{enumerate}
\item \textbf{Orbit} ($z$) --- \emph{coordinate precision} (\code{mpf\_t},
  hundreds to tens of thousands of bits):
  $z_{n+1} = z_n^2 + c$.
\item \textbf{First derivative} ($\partial z/\partial c$, written
  $\mathit{dzdc}$) --- \emph{derivative precision} (\code{mpf\_t}, a
  reduced precision $P_d \le P_c$; see \cref{subsec:ff-mixed-prec}):
  $\mathit{dzdc}_{n+1} = 2\,z_n \cdot \mathit{dzdc}_n + 1$.
\item \textbf{Second derivative} ($\partial^2 z/\partial c^2$, written
  $d2$) --- \emph{low precision} (\code{HDRFloat<double>}, ${\sim}53$-bit
  mantissa with wide exponent range):
  $d2_{n+1} = 2\bigl(\mathit{dzdc}_n^{\,2} + z_n \cdot d2_n\bigr)$.
\end{enumerate}

\noindent
Initial conditions: $z_0 = 0$, $\;\mathit{dzdc}_0 = 0$, $\;d2_0 = 0$.

The second derivative is deliberately kept at low precision because it
enters only two quantities---the Halley correction factor
(\cref{subsec:ff-halley}) and the error estimator
(\cref{eq:ff-err}).  Both consumers depend primarily on the
\emph{exponent} of $d2$, not its mantissa, so even ${\sim}24$~bits of
mantissa (as on the GPU path) suffice; see \cref{subsec:ff-halley}
for a detailed analysis.
The $z$ and $\mathit{dzdc}$ values are truncated to low-precision
\code{HDRFloat} before entering the $d2$ recurrence---the $d2$
computation never touches the high-precision \code{mpf\_t}
representation.

After $p$~iterations the outputs $z_p$, $\mathit{dzdc}_p$, and $d2_p$ are
returned to the outer loop.

\paragraph{Real/imaginary expansion.}
To implement the inner loop with ordinary (real-valued) arithmetic, write
\[
z = x + i\,y,\qquad
\mathit{dzdc} = dx + i\,dy,\qquad
d2 = d2x + i\,d2y,\qquad
c = c_x + i\,c_y.
\]
Each iteration step then consists of the following scalar operations.

\medskip\noindent
\textbf{Orbit}\quad ($z \leftarrow z^2 + c$):
\begin{equation}\label{eq:ff-orbit-ri}
\boxed{\begin{aligned}
x_{n+1} &= x_n^2 - y_n^2 + c_x, \\
y_{n+1} &= 2\,x_n\,y_n + c_y.
\end{aligned}}
\end{equation}

\noindent
\textbf{First derivative}\quad ($\mathit{dzdc} \leftarrow 2\,z\cdot\mathit{dzdc} + 1$):
\begin{equation}\label{eq:ff-dzdc-ri}
\boxed{\begin{aligned}
dx_{n+1} &= 2\,x_n \cdot dx_n - 2\,y_n \cdot dy_n + 1, \\
dy_{n+1} &= 2\,x_n \cdot dy_n + 2\,y_n \cdot dx_n.
\end{aligned}}
\end{equation}

\noindent
\textbf{Second derivative}\quad
($d2 \leftarrow 2\,(\mathit{dzdc}^{\,2} + z \cdot d2)$):
\begin{equation}\label{eq:ff-d2-ri}
\boxed{\begin{aligned}
d2x_{n+1} &= 2\,dx_n^2 - 2\,dy_n^2
              + 2\,x_n \cdot d2x_n - 2\,y_n \cdot d2y_n, \\
d2y_{n+1} &= 4\,dx_n \cdot dy_n
              + 2\,x_n \cdot d2y_n + 2\,y_n \cdot d2x_n.
\end{aligned}}
\end{equation}

\noindent
Each line is a sequence of multiplies, adds, and subtracts on real-valued
scalars---no complex arithmetic library is required.  Note that
\emph{within} a single iteration, the orbit update uses $x_n$,~$y_n$ (the
old values), the $\mathit{dzdc}$ update uses $x_n$,~$y_n$ and the old
$dx_n$,~$dy_n$, and the $d2$ update uses the old $x_n$,~$y_n$,
$dx_n$,~$dy_n$, $d2x_n$,~$d2y_n$.  The $d2$ recurrence must therefore be
evaluated \emph{before} the orbit and $\mathit{dzdc}$ values are
overwritten (or all old values must be saved).

\paragraph{Operation count.}
Not all operations are equally expensive.  The orbit and first-derivative
recurrences use high-precision \code{mpf\_t} arithmetic, where each
$N$-limb multiply costs $O(N \log N)$ (NTT-based) and each add/subtract
costs $O(N)$.  The second derivative uses low-precision
\code{HDRFloat} arithmetic (${\sim}53$-bit on the CPU path,
${\sim}24$-bit on the GPU path)---a handful of scalar floating-point
operations per iteration, negligible by comparison.

Counting real-valued scalar operations from
\cref{eq:ff-orbit-ri,eq:ff-dzdc-ri,eq:ff-d2-ri}, each iteration requires:

\begin{center}
\small
\begin{tabular}{llccc}
\hline
\textbf{Recurrence} & \textbf{Precision} &
  \textbf{Muls} & $\boldsymbol{\times 2}$ & \textbf{Add/subs} \\
\hline
Orbit ($z^2 + c$) & coord &
  3 & 1 & 3 \\
$\mathit{dzdc}$ ($2z \cdot \mathit{dzdc} + 1$) & deriv &
  4 & 2 & 3 \\
$d2$ ($2(\mathit{dzdc}^2 {+} z {\cdot} d2)$) & low-prec HDR &
  7 & 3 & 5 \\
\hline
\textbf{HP total} & &
  \textbf{7} & \textbf{3} & \textbf{6} \\
\hline
\end{tabular}
\end{center}

\noindent
``Muls'' are full $N \times N$ multiplications (\code{mpf\_mul}, $O(N \log
N)$ each)---the dominant cost.  ``$\times 2$'' denotes scaling by~two
(\code{mpf\_mul\_ui} on CPU, \code{addTwos} normalization on GPU), which is
$O(N)$ like an addition.  ``Add/subs'' includes all additions,
subtractions, and constant additions ($+1$, $+c_x$, $+c_y$), each $O(N)$.
The ``HP total'' row counts only the high-precision (coord and deriv)
operations; the $d2$ row contributes only cheap low-precision work.

\noindent
In the expanded equations each $\mathit{dzdc}$ and $d2$ term carries an
explicit factor of~2 (e.g.\ $2\,x \cdot dx$).  In the implementation
the~$\times 2$ is \emph{not} applied per-term: on the CPU the factor is
pre-applied to the operands ($2x$, $2y$ via \code{mpf\_mul\_ui}), and all
four derivative products share these two scaled operands.  On the GPU the
multiply kernel produces the raw products and the $\times 2$ is folded
into the normalization step (\code{addTwos}).  Either way the total is
4~full multiplies plus 2~$\times 2$ operations for $\mathit{dzdc}$.

\noindent
Concretely, the orbit contributes 3~multiplies at coordinate precision
($x^2$, $y^2$, $x \cdot y$) plus 3~add/subs ($x^2 - y^2$, $+c_x$,
$+c_y$) and one~$\times 2$ ($2 \cdot xy$).  The derivative contributes
4~multiplies at derivative precision ($x \cdot dx$, $y \cdot dy$,
$x \cdot dy$, $y \cdot dx$) plus 3~add/subs (one subtraction, one
addition, $+1$) and two~$\times 2$ (scaling $x$ and $y$ by~2 before
multiplying, or equivalently scaling the products after).

\noindent
The 7~high-precision multiplies are \emph{independent} of one another
within a single iteration and can all execute in parallel.  This independence is
exactly why the multithreaded CPU evaluator dispatches them to
7~spin-based worker threads (\cref{subsec:ff-threading}), and why the
GPU backend can batch them as simultaneous NTT launches
(\cref{subsec:ff-inner-gpu}).  The $d2$ computation runs on the main
thread (or a single GPU thread) concurrently with the high-precision
work, using truncated ${\sim}53$-bit copies of~$z$ and $\mathit{dzdc}$,
and completes long before the \code{mpf\_mul} workers finish.

\paragraph{The \code{zcoeff} product.}
In addition to $\partial z_p/\partial c$, the code tracks a multiplicative
coefficient
\begin{equation}\label{eq:ff-zcoeff}
w_n = \prod_{k=1}^{n-1} 2\,z_k,
\end{equation}
with the convention $w_0 = w_1 = 1$.  This product satisfies
$\partial z_n/\partial c = w_n + \text{lower-order terms}$ and is used later
to compute the intrinsic scale of the mini-brot (see
\cref{subsec:ff-intrinsic-radius}).  In the code, \code{zcoeff} is
initialized to~$1$ at $n=0$ (because $z_0 = 0$ would otherwise collapse the
product) and then multiplied by $2z$ at each subsequent step, matching the
recurrence for the leading factor of the derivative.

\subsubsection{GPU acceleration of the inner loop}
\label{subsec:ff-inner-gpu}

The inner loop's 7~high-precision multiplies per iteration are the
dominant cost of the entire Newton--Raphson procedure.  At deep zooms,
where the coordinate precision reaches thousands of 32-bit limbs, each
multiply is an $O(N \log N)$ NTT-based operation.  (The 6~additional
multiplies for the second derivative are low-precision \code{HDRFloat}
operations---${\sim}24$-bit on the GPU, ${\sim}53$-bit on the
CPU---and contribute negligible cost.)

A naive GPU approach would compute the 4~NR derivative products in a
\emph{separate} kernel launch from the 3~orbit products, doubling the
number of grid-wide synchronization barriers per iteration.  Because
barrier latency on a cooperative grid is largely independent of the
arithmetic performed between barriers, this would roughly halve the
effective throughput.

\FractalShark{}'s implementation avoids this penalty entirely: the
NR derivative products are \emph{interleaved} into the same multiply and
add kernel calls that already exist for the orbit-only reference orbit
computation, reusing the same barriers with zero additional
synchronization overhead.

\paragraph{Per-iteration structure.}
Each iteration of \code{ReferenceHelper}
(in \code{Kernel\-Hp\-Shark\-Reference\-Orbit\_cu.h})
proceeds in three phases:

\begin{enumerate}
\item \textbf{$d2$ accumulation (HDRFloat).}
  A single GPU thread (block~0, thread~0) computes the second derivative
  update (\cref{eq:ff-d2-ri}) in ${\sim}24$-bit \code{HDRFloat<float>}
  arithmetic (since \code{SharkParamsNR} uses \code{SubType=float}),
  using truncated copies of the current $z$ and
  $\mathit{dzdc}$.  This runs \emph{before} the multiply kernel and
  requires no synchronization: the $d2$ values are private to one thread
  and the multiply does not depend on them.

\item \textbf{Multiply kernel} (\code{MultiplyHelperNTTV2Separates}).
  Computes all 7~products in a single cooperative-grid launch:
  \begin{center}
  \small
  \begin{tabular}{lll}
  \hline
  \textbf{Product} & \textbf{Operands} & \textbf{Purpose} \\
  \hline
  XX & $x \times x$ & orbit: $x^2$ \\
  YY & $y \times y$ & orbit: $y^2$ \\
  XY & $x \times y$ & orbit: $x \cdot y$ \\
  W0 & $dx \times x$ & $\mathit{dzdc}$: $dx \cdot x$ \\
  W1 & $dy \times y$ & $\mathit{dzdc}$: $dy \cdot y$ \\
  W2 & $dx \times y$ & $\mathit{dzdc}$: $dx \cdot y$ \\
  W3 & $dy \times x$ & $\mathit{dzdc}$: $dy \cdot x$ \\
  \hline
  \end{tabular}
  \end{center}

\item \textbf{Add kernel} (\code{AddHelperSeparates}).
  Combines the 7~products into 4~results (2~orbit $+$ 2~derivative):
  \begin{center}
  \small
  \begin{tabular}{ll}
  \hline
  \textbf{Output} & \textbf{Computation} \\
  \hline
  $x_{n+1}$ & $\text{XX} - \text{YY} + c_x$ \\
  $y_{n+1}$ & $2 \cdot \text{XY} + c_y$ \\
  $dx_{n+1}$ & $2 \cdot \text{W0} - 2 \cdot \text{W1} + 1$ \\
  $dy_{n+1}$ & $2 \cdot \text{W2} + 2 \cdot \text{W3}$ \\
  \hline
  \end{tabular}
  \end{center}
\end{enumerate}

\paragraph{Barrier reuse in the multiply kernel.}
The NTT-based multiply pipeline
(\cref{sec:ntt-multiply}) requires grid-wide barriers between each
major phase: packing/twisting, forward NTT, pointwise multiply, inverse
NTT, unpacking, and normalization.  When NR is enabled, the forward and
inverse NTT passes use \code{FourWay} mode---processing
$x$, $y$, $dx$, $dy$ through the \emph{same} butterfly stages
with the \emph{same} internal barriers.  The pointwise multiply
phase computes all 7~products in a \emph{single} grid-strided loop with
one trailing barrier.  Unpacking and normalization are handled by
\code{NWay<7>} variants that process all 7~results in one pass.

The result: \textbf{the multiply kernel executes the same number of
\code{grid.sync()} barriers whether NR is enabled or not.}  The NR
products add only a small amount of pointwise Montgomery arithmetic
per frequency index---negligible compared to the barrier and memory
costs.

\paragraph{Barrier reuse in the add kernel.}
The add kernel computes the orbit results ($x^2 - y^2 + c_x$ and
$2xy + c_y$) via three-operand addition with magnitude comparison,
sign selection, carry propagation (parallel prefix scan), and
normalization.  When NR is enabled, the W0--W3 arrays are passed
alongside the orbit arrays into the \emph{same} function calls:
\code{Phase1\_ABC} handles both $\text{XX} - \text{YY} + c_x$ and
$\text{W0} - \text{W1} + 1$ in one pass;
\code{Phase1\_DE} handles both $2\text{XY} + c_y$ and
$\text{W2} + \text{W3}$ in one pass;
\code{CarryPropagation\_ABC\_PPv5} resolves carries for all
6~streams (3~orbit $+$ 3~NR) simultaneously.

Again, \textbf{no additional barriers are introduced.}  The NR arrays
are separate buffers processed in parallel with the orbit arrays
between the same synchronization points.

\paragraph{The \code{addTwos} mechanism.}
The factor of~2 in the expanded equations
(\cref{eq:ff-orbit-ri,eq:ff-dzdc-ri})---$2xy$,
$2\,x \cdot dx$, etc.---is \emph{not} applied as a separate
multiply-by-two pass.  Instead, the normalize phase of the multiply
kernel applies it via an \code{addTwos} exponent adjustment array:
\[
\code{addTwos[7]} = \{0,\; 0,\; 1,\; 1,\; 1,\; 1,\; 1\},
\]
where a value of~1 shifts the result exponent by one bit (equivalent to
$\times 2$).  The orbit products XX and YY are not scaled; XY and all
four W products receive the $\times 2$.  This layout matches the distributed
form of the equations: each derivative product appears with a factor
of~2 already absorbed.

\paragraph{Summary: zero-overhead NR derivatives.}
The orbit-only reference orbit kernel computes 3~multiplies and
2~add-results per iteration.  Enabling NR adds 4~more multiplies and
2~more add-results, but these are folded into the existing pipeline
without any new kernel launches or synchronization barriers.  The GPU
thread utilization improves (more arithmetic per barrier), and the
total barrier count per iteration---the dominant performance
factor---is \emph{unchanged}.

\medskip
\noindent
The outer Newton loop (typically ${\sim}10$--$30$ iterations) runs on
the host; the inner loop---$p$ iterations of high-precision
arithmetic per Newton step---runs on the GPU.  The host reads back
$z_p$, $\mathit{dzdc}_p$, and $d2_p$ after each $p$-iteration batch,
computes the Newton/Halley step, updates~$c$, and relaunches the inner
loop.  For large $p$ and $N$ the total arithmetic volume scales as
$O(p \cdot N \log N)$ per Newton step, making the GPU's throughput
advantage decisive.  The dispatch function
\code{DispatchNRByPrecision} rounds the coordinate precision to
the nearest supported \code{SharkParamsNR} limb count (256 through
524\,288) and invokes the corresponding GPU kernel.

\begin{figure}[!htbp]
\centering
\begin{tikzpicture}[
    >=Stealth,
    phase/.style={draw, rounded corners, minimum width=2.2cm, minimum height=0.6cm,
                  align=center, font=\scriptsize},
    orbit/.style={phase, fill=blue!12},
    nr/.style={phase, fill=orange!15},
    barrier/.style={draw, red!70!black, line width=1.5pt, dashed},
    barrlbl/.style={font=\tiny\bfseries, red!70!black},
    prod/.style={font=\tiny, align=center},
]
  % Phase labels across top
  \node[font=\small\bfseries] at (0, 2.7) {Pack/Twist};
  \node[font=\small\bfseries] at (2.8, 2.7) {Forward NTT};
  \node[font=\small\bfseries] at (5.6, 2.7) {Pointwise mul};
  \node[font=\small\bfseries] at (8.4, 2.7) {Inverse NTT};
  \node[font=\small\bfseries] at (11.2, 2.7) {Normalize};

  % Orbit row
  \node[orbit] (p1o) at (0, 1.5) {$x, y$};
  \node[orbit] (p2o) at (2.8, 1.5) {FourWay};
  \node[orbit] (p3o) at (5.6, 1.5) {XX, YY, XY};
  \node[orbit] (p4o) at (8.4, 1.5) {ThreeWay};
  \node[orbit] (p5o) at (11.2, 1.5) {7-way};

  % NR row
  \node[nr] (p1n) at (0, 0.5) {$dx, dy$};
  \node[nr] (p2n) at (2.8, 0.5) {(shared)};
  \node[nr] (p3n) at (5.6, 0.5) {W0--W3};
  \node[nr] (p4n) at (8.4, 0.5) {FourWay};
  \node[nr] (p5n) at (11.2, 0.5) {(shared)};

  % Arrows
  \foreach \i/\j in {p1o/p2o, p2o/p3o, p3o/p4o, p4o/p5o} \draw[->] (\i) -- (\j);
  \foreach \i/\j in {p1n/p2n, p2n/p3n, p3n/p4n, p4n/p5n} \draw[->] (\i) -- (\j);

  % Barriers
  \foreach \x in {1.4, 4.2, 7.0, 9.8} {
    \draw[barrier] (\x, 0.0) -- (\x, 2.2);
    \node[barrlbl] at (\x, -0.2) {sync};
  }

  % Row labels
  \node[font=\scriptsize, blue!70!black, rotate=90] at (-1.3, 1.5) {Orbit};
  \node[font=\scriptsize, orange!70!black, rotate=90] at (-1.3, 0.5) {NR deriv};

  % Add kernel below
  \draw[->] (p5o.south) -- ++(0, -0.5) coordinate (addstart);
  \draw[->] (p5n.south) -- ++(0, -0.5);

  \node[orbit, minimum width=4cm] (addo) at (9.0, -1.0)
    {$x^2{-}y^2{+}c_x$,\; $2xy{+}c_y$};
  \node[nr, minimum width=4cm] (addn) at (9.0, -1.8)
    {W0${-}$W1${+}$1,\; W2${+}$W3};

  \node[font=\small\bfseries] at (5.5, -1.4) {Add kernel};
  \draw[barrier] (7.0, -2.3) -- (7.0, -0.5);
  \node[barrlbl] at (7.0, -2.5) {sync (shared)};

  \draw[->] (addo) -- ++(2, 0) node[right, font=\scriptsize] {$x_{n+1}, y_{n+1}$};
  \draw[->] (addn) -- ++(2, 0) node[right, font=\scriptsize] {$dx_{n+1}, dy_{n+1}$};
\end{tikzpicture}
\caption[GPU multiply+add pipeline with shared barriers]{GPU multiply+add pipeline for one iteration.  The orbit
  products (blue, top) and NR derivative products (orange, bottom) flow
  through the \emph{same} NTT pipeline phases, sharing all grid-wide
  barriers (red dashed lines).  The add kernel similarly interleaves
  orbit and NR combines.  Enabling NR adds arithmetic between barriers
  but introduces no new synchronization points.}
\label{fig:ntt-pipeline}
\end{figure}

\subsection{Halley's method}
\label{subsec:ff-halley}

The second derivative $F'' = \partial^2 z_p/\partial c^2$ is computed during
every forward evaluation, even when the Halley step itself is not taken.  This
is because $F''$ serves two distinct purposes in the refinement loop:

\begin{enumerate}
\item \textbf{Halley correction.}  When the Halley gate (\cref{eq:ff-halley-gate})
  is satisfied, $F''$ enables a cubic-convergence step that can save expensive
  forward evaluations compared to plain Newton (see below).
\item \textbf{Error estimator.}  The stopping criterion
  (\cref{eq:ff-err}) uses $|F''|^2$ to estimate the residual after
  each step.  Without the second derivative, the code would have no reliable
  way to judge convergence short of comparing successive iterates---which
  requires an extra forward evaluation per check.
\end{enumerate}

Because $F''$ is carried in \code{HDRFloat} (${\sim}53$~bits of
mantissa on the CPU path, ${\sim}24$~bits on the GPU path; both with a
wide exponent), it adds only cheap low-precision arithmetic to
the forward evaluation: a few \code{HDRFloat} multiplies and additions per
iteration of the period loop, compared with the expensive \code{mpf\_t}
multiplies needed for $z$ and $\partial z/\partial c$.  The cost is therefore
negligible relative to the rest of the evaluation, making it worthwhile to
compute unconditionally.

Phase~B optionally upgrades from Newton to Halley iteration for cubic
convergence.  The Halley step for $F(c) = z_p(c)$ is
\begin{equation}\label{eq:ff-halley-step}
\Delta_H \;=\;
  \frac{2\,F(c)\,F'(c)}
       {2\bigl(F'(c)\bigr)^2 - F(c)\,F''(c)},
\qquad c \leftarrow c - \Delta_H,
\end{equation}
where $F'=\partial z_p/\partial c$ and $F''=\partial^2 z_p/\partial c^2$.

To evaluate \eqref{eq:ff-halley-step}, the denominator $D = 2(F')^2 - F\,F''$ involves a complex squaring and a
complex product.  Writing $F = f_x + i\,f_y$, $F' = f'_x + i\,f'_y$, and
$F'' = f''_x + i\,f''_y$:
\begin{align}
D_x &= 2\,(f'^2_x - f'^2_y) - (f_x\,f''_x - f_y\,f''_y), \\
D_y &= 2\,(2\,f'_x\,f'_y) - (f_x\,f''_y + f_y\,f''_x).
\end{align}
The numerator $N = 2\,F\,F'$ is a complex product, and the final step
$\Delta_H = N / D$ is a complex division as in \cref{eq:ff-newton-step}.

\paragraph{Second derivative recurrence.}
The second derivative $d2 = \partial^2 z/\partial c^2$ is computed by the
inner loop (\cref{subsubsec:ff-inner}) using the recurrence in
\cref{eq:ff-d2-ri}.  In the code this quantity is carried as a pair of
\code{HDRFloat} values (real and imaginary parts), which provide a
wide exponent range at low mantissa precision---\code{HDRFloat<double>}
(${\sim}53$-bit) on the CPU path and \code{HDRFloat<float>}
(${\sim}24$-bit) on the GPU path.  The coordinate-precision
quantities $z$ and $\partial z/\partial c$ are stored in \code{mpf\_t},
while $\partial^2 z/\partial c^2$ remains in HDR throughout the forward
evaluation.  When the Halley step is computed, the HDR second
derivative is promoted to \code{mpf\_t} at coordinate precision.

\paragraph{Halley gate.}
Halley's step is only beneficial when $F''$ is well-conditioned relative to
$F$ and $F'$.  The code uses the dimensionless ratio
\begin{equation}\label{eq:ff-halley-gate}
\rho^2 \;=\;
  \frac{|F|^2\,|F''|^2}{|F'|^4}
\end{equation}
and activates Halley only when $\rho^2 < 2^{-12}$.  When $\rho^2$ is
small the denominator $2(F')^2 - F\,F''$ is dominated by the $(F')^2$ term,
so the Halley correction is a small, well-determined perturbation of the
Newton step.  If the Halley denominator turns out to be exactly zero (or the
gate is not satisfied), the code falls back to a plain Newton step for that
iteration.

\paragraph{Is Halley worthwhile given low-precision $F''$?}
The second derivative $F''$ is accumulated in \code{HDRFloat}
throughout the forward evaluation, providing only $b$~bits of
mantissa (with a wide exponent range): $b \approx 53$ on the CPU path
(\code{HDRFloat<double>}) and $b \approx 24$ on the GPU path
(\code{HDRFloat<float>}, since \code{SharkParamsNR} uses
\code{SubType=float}).  Meanwhile $F$ and $F'$ are carried in
\code{mpf\_t} at hundreds or thousands of bits.  This precision mismatch
limits the benefit of the Halley correction.

The Halley step can be decomposed as
\begin{equation}\label{eq:ff-halley-decomp}
\Delta_H \;=\; \Delta_N \cdot \frac{1}{1 - \frac{F \cdot F''}{2(F')^2}}
\;=\; \Delta_N \cdot \frac{1}{1 - \tfrac{1}{2}\rho\, e^{i\theta}},
\end{equation}
where $\Delta_N = F/F'$ is the Newton step, $\rho = \sqrt{\rho^2} =
|F|\cdot|F''|/|F'|^2$ is the first-power ratio (the gate criterion checks
$\rho^2$; the correction factor has magnitude $\rho/2$), and $e^{i\theta}$
is the complex phase of $F \cdot F'' / (F')^2$.

\paragraph{Why it works: error analysis of the correction factor.}
Starting from \eqref{eq:ff-halley-decomp}, write the true second derivative as $F''$ and the $b$-bit approximation as
$\hat{F}'' = F''(1 + \epsilon)$ where $|\epsilon| \lesssim 2^{-b}$.  The
Halley step with the approximate $\hat{F}''$ is
\[
\hat{\Delta}_H \;=\; \Delta_N \cdot
  \frac{1}{1 - \frac{F \cdot \hat{F}''}{2(F')^2}}.
\]
Near a simple root, the error $e_k = c_k - c^*$ satisfies $F \approx F' e_k$,
so the correction ratio
\[
\frac{F \cdot F''}{2(F')^2} \;\approx\;
  \frac{e_k \cdot F''}{2 F'} \;=\; O(e_k).
\]
This term shrinks as $c_k$ converges.  The \emph{error} introduced by using
$\hat{F}''$ instead of $F''$ modifies this ratio by
\[
\frac{F \cdot \hat{F}''}{2(F')^2}
  - \frac{F \cdot F''}{2(F')^2}
  \;=\; \frac{F \cdot F'' \epsilon}{2(F')^2}
  \;=\; O(e_k \cdot \epsilon).
\]
This quantity is the product of the current error and the $b$-bit truncation noise.
The resulting Halley step error, relative to the true Halley step, is therefore
$O(e_k \cdot \epsilon)$, producing a next error of
\[
e_{k+1}^{\text{(approx Halley)}} \;=\;
  \underbrace{O(e_k^3)}_{\text{true Halley}}
  \;+\;
  \underbrace{O(e_k^2 \cdot \epsilon)}_{\text{truncation of } F''}.
\]
(The extra factor of $e_k$ in the second term comes from the fact that
$\Delta_N$ itself is $O(e_k)$, and the relative perturbation of the
correction factor is $O(e_k \cdot \epsilon)$.)

Two regimes emerge:
\begin{itemize}
\item \textbf{Early convergence} ($|e_k| \gg 2^{-b}$): the $O(e_k^3)$ term
  dominates, and the method converges cubically---identical to full-precision
  Halley.  The $b$-bit noise is invisible.
\item \textbf{Late convergence} ($|e_k| \ll 2^{-b}$): the
  $O(e_k^2 \cdot \epsilon)$ term dominates.  Since $\epsilon$ is a fixed
  ${\sim}2^{-b}$, the effective recurrence becomes $|e_{k+1}| \approx
  C \cdot e_k^2 \cdot 2^{-b}$, which is quadratic convergence (like Newton)
  but with an extra constant factor of $2^{-b}$.  In terms of correct bits,
  this gives ${\sim}2n + b$ bits per step---still better than Newton's
  ${\sim}2n$, but no longer cubic.
\end{itemize}

The crossover occurs when $|e_k|^3 \approx |e_k|^2 \cdot 2^{-b}$,
i.e.\ when $|e_k| \approx 2^{-b}$, or equivalently when the iterate has
${\sim}b$ correct bits.

\paragraph{Empirical confirmation ($b = 53$).}
Numerical experiments on the period-3 nucleus confirm this analysis
for the CPU path ($b = 53$, \code{HDRFloat<double>}):

\begin{center}
\begin{tabular}{crrr}
\hline
\textbf{Step} & \textbf{Newton} & \textbf{Halley (53-bit $F''$)}
              & \textbf{Halley (full $F''$)} \\
\hline
0 &  13.8 &  20.7 &  20.7 \\
1 &  26.5 &  60.4 &  60.4 \\
2 &  51.9 & 173.8 & 179.6 \\
3 & 102.6 & 400.6 & 536.9 \\
4 & 204.1 & 854.2 & 1608.9 \\
5 & 407.0 & 1761.4 & --- \\
6 & 812.8 & --- & --- \\
7 & 1624.4 & --- & --- \\
\hline
\end{tabular}
\end{center}

(Correct bits of the period-3 nucleus, starting from the same initial guess.)
Steps~0--1 show cubic convergence for both Halley variants (ratio
${\approx}\,3$). At step~2, the 53-bit Halley begins to diverge from the
full-precision version (173.8 vs.\ 179.6), consistent with the crossover at
${\sim}53$ correct bits.  By step~4, the 53-bit Halley's ratio has dropped to
${\sim}2.1$, approaching quadratic.

\paragraph{Net benefit.}
On the period-3 example, the low-precision Halley reaches 1761~correct bits
in 6~steps versus Newton's 8~steps for 1624~bits---saving two full forward
evaluations.  This saving comes from the early cubic burst: by the time the
method degenerates to quadratic, it already has a ${\sim}3\times$ lead in
correct bits over Newton, and this lead persists through the later quadratic
phase.  Such savings are most pronounced when Phase~B starts below the
${\sim}b$-bit crossover.  At already-accurate starts, the large cubic lead is
absent, so the retained benefit is narrower: representative runs show only a
0--1 outer-iteration advantage in this late regime, while $F''$ still remains
useful for the Halley gate and error estimator.  For large period~$p$, even
that one saved forward evaluation still avoids $p$~expensive
high-precision multiplies.

\paragraph{Dependence on initial-guess quality.}
Whether this step-count saving is significant in practice depends on how many
correct bits the Phase~B starting point already has.  Phase~A produces an
initial candidate from the renderer's pixel grid and a few Newton steps in
the native float type~$T$.  The pixel resolution contributes approximately
$\lfloor\log_2 Z\rfloor - \lceil\log_2 W\rceil$ correct bits, where $Z$ is
the zoom factor and $W$ the screen width in pixels.  At shallow zoom
($Z \sim 10^{4}$--$10^{10}$), the initial guess has roughly 14--30~correct
bits, well below the ${\sim}b$-bit crossover, and the cubic burst is fully
observable during Phase~B; this is the regime in which low-precision Halley
delivers most of its practical iteration savings.  At deep zoom ($Z \gg 2^b$),
pixel resolution alone provides hundreds of correct bits, placing the iterate
well past the crossover before Phase~B begins.  In that regime the cubic
convergence advantage is absent; the retained value of low-precision Halley is
limited to the smaller late-regime improvement over Newton together with the
information carried by $|F''|$ for the Halley gate and error estimator.
Representative period~16045 runs still show the dependence.  After accounting
for the ${\sim}88$-bit offset between actual and normalized bits
(\cref{eq:ff-rho2-offset}), a 100-bit start retains about a one-iteration
advantage to representative targets, whereas a 150-bit start is already in the
narrowed 0--1-iteration envelope.  This result is consistent with the implementation
and with partial production spot checks, but it is not yet a full validation of
every production regime.

\paragraph{Interpreting the diagnostic metric.}
The diagnostic ratio $\rho^2$ (\cref{eq:ff-halley-gate}) measures
convergence in a normalized sense rather than counting correct bits of
the candidate~$c$.  Near a simple root, $F \approx F' e_k$ where
$e_k = c_k - c^*$, so
\begin{equation}\label{eq:ff-rho2-offset}
\rho^2 \;\approx\; |e_k|^2 \cdot \frac{|F''|^2}{|F'|^2},
\qquad
  \text{normalized bits}
  \;=\; -\tfrac{1}{2}\log_2 \rho^2
  \;=\; \text{correct bits} - \log_2 \frac{|F''|}{|F'|}.
\end{equation}
For the period-3 nucleus, $|F''/F'| \approx 2^{2.2}$ and the offset is
negligible.  At period~16045, $|F''/F'| \approx 2^{88}$, so the
$\rho^2$-derived metric understates actual correct bits by approximately
88~bits.  At extreme periods ($p \sim 10^8$), the offset grows further,
which can mask the early cubic burst in the diagnostic: the Phase~A
initial guess may already have enough actual correct bits to place the
iterate past the ${\sim}b$-bit crossover, even though the normalized
metric appears low.

\paragraph{GPU variant ($b = 24$).}
The GPU NR path (\code{SharkParamsNR}) uses \code{SubType=float},
so $d2$ is accumulated in \code{HDRFloat<float>} with $b \approx 24$.
The crossover from cubic to quadratic convergence therefore occurs
at ${\sim}24$~correct bits rather than ${\sim}53$, and the late regime
gives ${\sim}2n + 24$ bits per step instead of ${\sim}2n + 53$.
Representative fp32-like sweeps suggest that the remaining practical
step-count benefit is likely even more concentrated toward comparatively coarse
starts than on the CPU-like path.  Relative to the 53-bit CPU-like variant,
the 24-bit path can forfeit roughly one outer iteration in some targets, but
at already-accurate starts both variants often collapse to the same narrow
0--1-iteration envelope relative to Newton.  This trend is consistent with the
implementation and with partial production spot checks, but it is not yet a
full production validation.

Using \code{double} or \code{dblflt} for the GPU $d2$ accumulation
would shift the crossover to ${\sim}53$ or ${\sim}48$ correct bits,
saving that one iteration.
However, the $d2$ accumulation runs on a single GPU thread
(block~0, thread~0) during the multiply phase, and upgrading it would
increase register pressure and arithmetic latency for a gain of at most
one saved forward evaluation.  This tradeoff is not currently considered
worthwhile.

The error estimator (\cref{eq:ff-err}) and Halley gate
(\cref{eq:ff-halley-gate}) are unaffected by this choice: both depend
only on the \emph{exponent} of $|F''|$, which the \code{HDRFloat}
representation tracks exactly regardless of mantissa precision.

\subsection{Mixed-precision refinement}
\label{subsec:ff-mixed-prec}

The Phase~B polisher (\code{RefinePeriodicPoint}) operates entirely in GMP
\code{mpf\_t} arithmetic but uses two distinct precision levels to balance cost
and accuracy, following the strategy used in Imagina.

\paragraph{Coordinate vs.\ derivative precision.}
The candidate coordinate $c$ and the orbit value $z_p$ are maintained at full
\emph{coordinate precision} $P_c$ (typically hundreds to tens of thousands of
bits, determined by the zoom depth).  The derivative
$\partial z_p/\partial c$ is maintained at a reduced \emph{derivative
precision} $P_d$ chosen by
\begin{equation}\label{eq:ff-deriv-prec}
P_d = \operatorname{clamp}\!\left(
  \frac{-E_{\mathrm{scale}} + 32}{4},\;
  256,\;
  P_c + E_{\max}^{|c|} + 32
\right),
\end{equation}
where $E_{\mathrm{scale}}$ is the base-2 exponent of the scale factor
$\text{Scale} \approx 1/|w_p \cdot \partial z_p/\partial c|$ and
$E_{\max}^{|c|}$ is the exponent of $\max(|\Re c|,|\Im c|)$.
The rationale is that the derivative need only be accurate enough to produce
a Newton/Halley step whose relative error is smaller than the current
residual; carrying it at full coordinate precision would roughly double the
cost of each forward evaluation without improving the convergence rate.

\paragraph{Error estimator and stopping criterion.}
After each step $\Delta c$, the code computes an error proxy entirely in
\code{HDRFloat<double>}:
\begin{equation}\label{eq:ff-err}
\mathrm{err} \;=\;
  \frac{|\Delta c|^4 \cdot |F''|^2}{|F'|^2}.
\end{equation}
Near a simple root the classical Newton error bound gives
$|e_{k+1}| \approx |F''|/(2|F'|) \cdot |\Delta c|^2$,
so $\mathrm{err} \approx 4\,|e_{k+1}|^2$.
Iteration stops when
\begin{equation}\label{eq:ff-stop}
-\lfloor \log_2(\mathrm{err}) \rfloor \;\ge\; 2\,P_c,
\end{equation}
meaning $\mathrm{err} < 2^{-2P_c}$ and hence
$|e_{k+1}| \lesssim 2^{-P_c}$---the next iterate has at least $P_c$~correct
bits.

\paragraph{Final correction and accept/reject.}
After the main loop converges (or exhausts the iteration budget), a single
additional Newton step is applied as a final correction.  The refined $c$ is
then checked against the initial seed $c_0$: if
$|c - c_0|^2 > R^2$ (the squared search radius, carried in \code{mpf\_t}),
the result is rejected and $c$ is reset to~$c_0$.  This guard prevents the
iteration from wandering to a distant nucleus outside the user's region of
interest.

\subsection{Intrinsic radius and scale}
\label{subsec:ff-intrinsic-radius}

Once the nucleus $c^*$ and its period~$p$ are known, the local magnification
of the mini-brot copy is characterized by
\begin{equation}\label{eq:ff-scale}
\mathrm{Scale} \;=\;
  \frac{1}{\bigl|w_p \cdot \partial z_p/\partial c\bigr|},
\end{equation}
where $w_p = \prod_{k=1}^{p-1} 2\,z_k$ is the \code{zcoeff} product
\eqref{eq:ff-zcoeff} evaluated at $c^*$. The code reports an \emph{intrinsic radius}
\begin{equation}\label{eq:ff-intrinsic-radius}
r_{\mathrm{int}} = 4 \cdot \mathrm{Scale}
  = \frac{4}{\bigl|w_p \cdot \partial z_p/\partial c\bigr|}.
\end{equation}
Geometrically, $r_{\mathrm{int}}$ estimates the diameter of the
period-$p$ hyperbolic component in the parameter plane.  It is stored
as a \code{HighPrecision} value in \code{FeatureSummary} and can be used by
the renderer to set an appropriate zoom level or to determine the additional
bits of precision needed to render the component's interior.

\subsection{Checkpoint and resume}
\label{subsec:ff-checkpoint}

High-precision Newton--Raphson refinement can require minutes to hours
at extreme periods, because each outer step evaluates the critical orbit
for $p$~iterations at the current working precision.  To avoid losing
progress when the user aborts (\code{Ctrl+Alt}), closes the
application, or encounters an error, \FractalShark{} writes periodic
checkpoints to a file (\code{nr\_checkpoint.txt}) and resumes from the
most recent checkpoint on the next invocation.

\paragraph{Checkpoint contents.}
Each checkpoint records the current candidate~$c$ (real and imaginary
parts as decimal strings at full working precision), the detected
period~$p$, the MPIR precision in bits, the outer Newton iteration
index, and the inner-loop iteration count within the current $p$-step
evaluation.  Storing the inner-loop position allows resumption
\emph{mid-evaluation}, not only at outer-step boundaries.

\paragraph{Asynchronous writer.}
Converting high-precision \code{mpf\_t} values to decimal strings via
\code{mpf\_get\_str} and writing them to disk can take non-trivial
time at high precision.  A dedicated \code{CheckpointWriter} background
thread accepts \code{CheckpointSnapshot} objects (which own independent
copies of the \code{mpf\_t} data) and performs the serialization and
file I/O without blocking the Newton iteration.  The snapshot is
written to a temporary file and atomically renamed to the checkpoint
path, ensuring that the on-disk checkpoint is always in a consistent
state.

\paragraph{Save points.}
Checkpoints are written at two granularities:
\begin{itemize}
  \item After each outer Newton step completes (precision doubling
        boundary).
  \item Periodically during the inner-loop evaluation, via a progress
        callback that triggers whenever a configurable number of
        inner iterations have elapsed.
\end{itemize}
Both paths route through the same asynchronous writer, so at most one
background write is in flight at a time.

\paragraph{Resume.}
On entry to \code{RefinePeriodicPoint}, the code checks for an existing
checkpoint file.  If the file exists and its stored period and precision
match the current request, the saved candidate~$c$ and iteration state
are loaded and the refinement loop skips ahead to the checkpointed
position.  If the period or precision do not match (indicating a
different feature or a re-run at different settings), the checkpoint is
ignored and refinement starts from scratch.

\paragraph{Abort safety.}
When the \code{AbortMonitor} sets the stop flag, the refinement loop
exits early \emph{after} writing a final checkpoint.  Progress is
therefore preserved: the user can later invoke \emph{Resume NR
Refinement} from the Feature Finder menu to continue from where the
computation was interrupted.

\subsection{Threading model}
\label{subsec:ff-threading}

\paragraph{Multithreaded CPU implementation.}
The dominant cost of the Feature Finder is the forward evaluation
\code{EvaluateCriticalOrbitAndDerivs}, which performs the inner loop
described in \cref{subsubsec:ff-inner} at high precision using GMP
\code{mpf\_t} arithmetic.  Because the 7~high-precision multiplies within
each iteration are independent (\cref{subsec:ff-inner-gpu}),
\FractalShark{} dispatches them to a pool of seven spin-based worker
threads (\code{NWORK~=~7})---one worker per multiply.
Synchronization uses cache-line-aligned (\code{alignas(64)}) generation
counters with atomic acquire/release semantics---no mutexes or condition
variables---so the main thread and workers coordinate with minimal latency.
The main thread composes the results after each batch completes and
advances to the next iteration.

A single-threaded fallback (\code{EvaluateCriticalOrbitAndDerivsST}) with
identical logic is also provided.

The Newton outer loop (\cref{subsubsec:ff-outer}) in
\code{FindPeriodicPoint\_Common} is inherently sequential: each Newton step
$c \leftarrow c - z_p / (\partial z_p / \partial c)$ depends on the
preceding evaluation's output, so successive steps cannot overlap.

GPU acceleration of the inner loop via the \code{HpSharkFloat} backend
is described in \cref{subsec:ff-inner-gpu}; this path is fully
implemented and dispatches through
\code{Dispatch\-NR\-By\-Precision} in \code{Feature\-Finder.cpp}.

With the core iteration algorithms and precision infrastructure now
described, the next chapter turns to the post-iteration processing
pipeline---antialiasing, color mapping, and rendering orchestration---that
transforms raw escape-time counts into the final displayed image.

